\subsection{Directional PAC--Bayes complexity and optimal flag variances}
The previous Gaussian bound uses one curvature constant for all directions.
An anisotropic envelope instead charges perturbations according to the
curvature of each flag block. We optimize the variance within every block,
then examine the effect of refining the flag and of tracking its oracle
inexactly. These are comparisons of posterior bounds under the stated
evaluation-loss assumptions.

\begin{assumption}[Anisotropic curvature envelope]\label{ass:directional-curvature}
For the bounded evaluation loss $0\le\ell(w,z)\le1$, assume a sample-independent
matrix $H_0\succeq0$ such that, globally in $w,z$,
\begin{equation}
 -H_0\preceq\nabla_w^2\ell(w,z)\preceq H_0.
\end{equation}
For $\mathcal P_j:=I_m\otimes\Pi_j(U)$, set
\begin{gather}
 d_j:=mq_j,\quad C_j(U):=\tr(\mathcal P_jH_0),\quad
 \overline h_j(U):=C_j(U)/d_j,\quad s_j:=\sigma_j^2>0,\\
 \Sigma_U:=\sum_j s_j\Pi_j(U),\quad
 \widetilde\Sigma_U:=I_m\otimes\Sigma_U=\sum_j s_j\mathcal P_j,\quad
 c_{n,\delta}:=\log(\Xi_n/\delta).
\end{gather}
The curvature assumption concerns the bounded evaluation loss; it does not
assert that unbounded training cross-entropy lies in $[0,1]$.
\end{assumption}
\begin{theorem}[Optimized directional posterior]\label{thm:directional-pac}
Under Assumption~\ref{ass:directional-curvature}, fix
$P=\mathcal N(0,\sigma_0^2I_N)$ independently of the sample. With probability
at least $1-\delta$, simultaneously for every learned $W,U$ and every $\lambda>0$,
\begin{align}
 F_\mathcal D(W)
 &\le F_S(W)+\frac{\|W\|_F^2}{2\lambda\sigma_0^2}
 +\frac{c_{n,\delta}}\lambda+\frac\lambda{8n}
 +\mathcal C_\lambda(U),\\
 \mathcal C_\lambda(U)
 &:=\frac1{2\lambda}\sum_{j=1}^{K+1}d_j
                 \log(1+2\lambda\sigma_0^2\overline h_j(U)).
 \label{eq:directional-complexity}
\end{align}
The optimal positive posterior variances in this bound are exactly
\begin{equation}
 s_j^\star(U,\lambda)
 =\frac{\sigma_0^2}{1+2\lambda\sigma_0^2\overline h_j(U)}.
\end{equation}
\end{theorem}
\begin{proof}
For $E\sim\mathcal N(0,\widetilde\Sigma_U)$, Taylor's formula gives
\begin{align}
 |\ell(W+E,z)-\ell(W,z)-\langle\nabla\ell(W,z),E\rangle|
 &\le\tfrac12 E^\top H_0E,\\
 |L(Q)-F_\mathcal D(W)|,\quad|\widehat L_S(Q)-F_S(W)|
 &\le\tfrac12\tr(H_0\widetilde\Sigma_U)=\tfrac12\sum_j C_js_j.
\end{align}
On the single event of Theorem~\ref{thm:rigid-pac},
\begin{align}
 F_\mathcal D(W)
 &\le F_S(W)+\sum_jC_js_j+
 \sqrt{\frac{\|W\|_F^2/(2\sigma_0^2)
   +\frac12\sum_jd_j[s_j/\sigma_0^2-1-\log(s_j/\sigma_0^2)]
   +c_{n,\delta}}{2n}},\\
 \sqrt{a/(2n)}&\le a/\lambda+\lambda/(8n)\quad(a\ge0,\lambda>0).
\end{align}
For each $j$, define the strictly convex scalar objective
\begin{gather}
 J_j(s):=C_js+\frac{d_j}{2\lambda}
                       [s/\sigma_0^2-1-\log(s/\sigma_0^2)],\\
 J_j'(s)=C_j+\frac{d_j}{2\lambda}(\sigma_0^{-2}-s^{-1}),\quad
 J_j''(s)=\frac{d_j}{2\lambda s^2}>0,\\
 J_j'(s_j^\star)=0,\quad
 J_j(s_j^\star)=\frac{d_j}{2\lambda}
                                  \log(1+2\lambda\sigma_0^2C_j/d_j).
\end{gather}
The event is uniform over all posteriors; optimization over data-dependent
$U,s_j,\lambda$ therefore requires no additional union bound.
\end{proof}
The optimal variance is smaller in blocks with greater average curvature:
perturbations there have a larger loss cost. The logarithmic complexity results
from balancing that cost against the KL penalty for shrinking variance.
Here $\lambda$ is a bound-optimization parameter; it is separate from the
optimizer's learning rate and from weight-decay notation used earlier.
The closed form optimizes this analytic envelope; an evaluated certificate
can be tighter by using the inverse-KL bound before its scalar relaxation.

\begin{proposition}[A strict benefit of posterior flag refinement]
\label{prop:directional-refinement}
Suppose a block $\mathcal P$ of dimension $d$ is refined into orthogonal blocks
$\mathcal P_1,\ldots,\mathcal P_J$, with
\begin{gather}
 \mathcal P=\sum_{j=1}^J\mathcal P_j,\quad d=\sum_jd_j,\quad
 h_j:=\tr(\mathcal P_jH_0)/d_j,\quad \overline h:=\sum_j(d_j/d)h_j.
\end{gather}
Then
\begin{equation}
 \sum_j\frac{d_j}{2\lambda}\log(1+2\lambda\sigma_0^2h_j)
 \le\frac d{2\lambda}\log(1+2\lambda\sigma_0^2\overline h).
\end{equation}
For $d_j>0$, equality holds if and only if $h_1=\cdots=h_J$. Thus refining a
fixed flag can strictly improve the optimized posterior bound when the block
curvatures differ. This concerns the posterior family, not an ordering of the
trained optimizers' test errors.
More quantitatively, let $h_-:=\min_jh_j$, $h_+:=\max_jh_j$,
$V_h:=\sum_j(d_j/d)(h_j-\overline h)^2$, and let $\mathcal G_\lambda$ be
the coarse cost minus the refined cost. Then
\begin{equation}
 \frac{d\lambda\sigma_0^4 V_h}{(1+2\lambda\sigma_0^2h_+)^2}
 \le\mathcal G_\lambda\le
 \frac{d\lambda\sigma_0^4 V_h}{(1+2\lambda\sigma_0^2h_-)^2}.
\end{equation}
The two sides agree to leading order as the block curvatures coalesce.
\end{proposition}
\begin{proof}
\begin{gather}
 f(h):=\log(1+2\lambda\sigma_0^2h),\qquad
 f''(h)=-\frac{(2\lambda\sigma_0^2)^2}{(1+2\lambda\sigma_0^2h)^2}<0,\\
 \sum_j(d_j/d)f(h_j)\le f(\sum_j(d_j/d)h_j).
\end{gather}
For the quantitative statement apply Taylor's formula at $\overline h$.
The weighted linear term cancels, and $-f''$ lies between its values at
$h_+$ and $h_-$. Multiplication by $d/(2\lambda)$ gives the two-sided
strong-concavity remainder. Continuity of $f''$ proves the local sharpness.
\end{proof}
Refinement permits different variances where a coarser posterior required a
shared one. Strict concavity makes this flexibility useful precisely when the
subblocks have unequal average curvature. Whether a covariance-trained flag
discovers such directions is a separate alignment question. The next result
controls the change in complexity caused by tracking a chosen oracle
inexactly, without assuming that oracle minimizes the curvature cost.

\begin{proposition}[Tracking error in the directional complexity]
\label{prop:directional-tracking}
Partition $H_0\in\RR^{md\times md}$ into $d\times d$ row-coordinate blocks
$H_{ab}$, and define $A_0:=\sum_{a=1}^m H_{aa}\succeq0$. Then
Set $s_A:=\lambda_{\max}(A_0)-\lambda_{\min}(A_0)$. Then
\begin{gather}
 C_j(U)=\tr(\Pi_j(U)A_0),\\
 |\mathcal C_\lambda(U)-\mathcal C_\lambda(V)|
 \le\sigma_0^2 s_A\sqrt d\,d_P(U,V).
\end{gather}
In the regular oracle chart, let $\widehat x_t^\star$ and $x_t^\star$ be the
oracles for $\widehat C_t$ and $C_t$, respectively, and $U_t=\psi(x_t)$. Then
\begin{align}
 |\mathcal C_\lambda(U_t)-\mathcal C_\lambda(\psi(x_t^\star))|
 &\le\sigma_0^2 s_A\sqrt d\,J_P
 \left(\|x_t-\widehat x_t^\star\|+
     \frac{2J_\psi}{\mu_\phi}\sum_k\omega_k\sqrt{r_k}\epsilon_{k,t}\right).
\end{align}
\end{proposition}
\begin{proof}
Centering by $c_A:=(\lambda_{\max}(A_0)+\lambda_{\min}(A_0))/2$ costs
nothing because every projector difference is traceless. The trace/nuclear
norm duality and the exact overlap identity give
For $g_j(c):=d_j(2\lambda)^{-1}\log(1+2\lambda\sigma_0^2c/d_j)$,
\begin{align}
 0\le g_j'(c)&=\frac{\sigma_0^2}{1+2\lambda\sigma_0^2c/d_j}\le\sigma_0^2,\\
 |C_j(U)-C_j(V)|
 &\le\|A_0-c_AI\|_{\op}\|\Pi_j(U)-\Pi_j(V)\|_*\\
 &\le\tfrac{s_A}{2}\sqrt{2q_j}\|\Pi_j(U)-\Pi_j(V)\|_F,\\
 \sum_{j=1}^{K+1}\|\Pi_j(U)-\Pi_j(V)\|_F^2
 &\le2d_P(U,V)^2,\\
 \sum_j\sqrt{2q_j}\|\Pi_j(U)-\Pi_j(V)\|_F
 &\le\sqrt{2\sum_jq_j}\sqrt{2d_P(U,V)^2}=2\sqrt d\,d_P(U,V).
\end{align}
Apply the triangle inequality and Lemma~\ref{lem:tf-oracle} to the two oracle branches.
In particular $A_0=cI$ gives zero orientation cost exactly. The bound depends
on spectral spread, removing an isotropic component that cannot distinguish flags.
\end{proof}
\begin{proposition}[A sample-independent anisotropic flag prior]
\label{prop:directional-prior}
Let a fixed prior flag $V$ have projectors $\Pi_l^0$ and ranks $q_l^0$, and let
\begin{gather}
 \Lambda_0:=\sum_l\tau_l^2\Pi_l^0\succ0,\quad
 P_V:=\mathcal N(0,I_m\otimes\Lambda_0),\quad
 Q_U:=\mathcal N(\operatorname{vec}W,I_m\otimes\Sigma_U),\\
 T_{lj}:=\tr(\Pi_l^0\Pi_j(U)),\quad T_{lj}\ge0,\quad
 \sum_jT_{lj}=q_l^0,\quad\sum_lT_{lj}=q_j.
\end{gather}
Then the exact orientation-dependent cost is
\begin{align}
 \operatorname{KL}(Q_U\|P_V)
 &=\tfrac12\tr(W\Lambda_0^{-1}W^\top)
 +\tfrac m2\left[\sum_{l,j}\frac{\sigma_j^2}{\tau_l^2}T_{lj}-d
 +\sum_lq_l^0\log\tau_l^2-\sum_jq_j\log\sigma_j^2\right].
\end{align}
A data-dependent $V$ is admissible if it is learned from an independent pilot
sample $S_0$ and the bound is evaluated on a disjoint iid sample $S_1$ of size $n_1$:
\begin{gather}
 S_0\perp S_1,\quad V=V(S_0),\\
 \mathcal E(S_0):=\left\{\begin{aligned}
 &\forall Q\ll P_{V(S_0)}:\\
 &\operatorname{kl}(\widehat L_{S_1}(Q)\|L(Q))
 \le\frac{\operatorname{KL}(Q\|P_{V(S_0)})+\log(\Xi_{n_1}/\delta)}{n_1}
 \end{aligned}\right\},\\
 \PP_{S_1\mid S_0}(\mathcal E(S_0))\ge1-\delta.
\end{gather}
The posterior may use both samples; the empirical term in this bound uses $S_1$.
\end{proposition}
\begin{proof}
\begin{align}
 \tr(\Lambda_0^{-1}\Sigma_U)
 &=\sum_{l,j}\tau_l^{-2}\sigma_j^2\tr(\Pi_l^0\Pi_j),\\
 T_{lj}&=\tr(\Pi_j\Pi_l^0\Pi_j)=\|\Pi_l^0\Pi_j\|_F^2\ge0,\\
 \log\det\Lambda_0-\log\det\Sigma_U
 &=\sum_lq_l^0\log\tau_l^2-\sum_jq_j\log\sigma_j^2.
\end{align}
Substitution in the Gaussian density-ratio identity proves the first assertion.
Conditioning on $S_0$ makes the prior independent of $S_1$; apply
Theorem~\ref{thm:rigid-pac} conditionally, then integrate over $S_0$.
\end{proof}
